% TEMPLATE-MILC-v3.tex - plantilla maestra UNICA de las guias MILC v3.
%
% La usan las tres familias de guias, cada una con su driver:
%   · Grados 6-11  -> scripts/build-guias-g11.py
%   · Bebras       -> scripts/build-guias-bebras.py
%   · Semillero    -> scripts/build-guias-semillero.py
%
% Antes habia tres copias casi identicas (~400 lineas iguales, 6 distintas):
% cada mejora habia que aplicarla tres veces, y en la practica no se hacia
% (el motor de recursos vivia solo en dos de ellas). Lo que varia entre
% familias esta parametrizado en 6 placeholders que cada driver llena
% (se nombran aqui SIN su sintaxis real: el builder reemplaza por texto plano,
% tambien dentro de los comentarios, y un valor multilinea rompe el preambulo):
%
%   HEADER_IZQ        encabezado izquierdo de pagina
%   HEADER_DER        encabezado derecho de pagina
%   PORTADA_ETIQUETA  rotulo sobre el numero grande (GRADO/MOMENTO/MODULO)
%   PORTADA_SERIE     linea de serie de la portada
%   PORTADA_ID        identificador de la guia en la portada
%   PORTADA_CIRCULO   el circulito de grado (°), o vacio si no aplica
%
% El resto de claves las documenta content/guias/_SCHEMA.md. El script valida
% que no queden placeholders sin reemplazar antes de escribir el archivo final.

\documentclass[11pt,letterpaper]{article}
\usepackage[margin=1.75cm]{geometry}
\usepackage[table]{xcolor}
\usepackage{fontspec}
\usepackage[spanish]{babel}
\usepackage{tikz}
% Librerías TikZ para los diagramas de las guías (bloque `recursos:`):
%  · babel  → babel en español vuelve activos caracteres como «>», y TikZ los usa
%             en sus opciones (>=stealth, ->). Sin esta librería, cualquier
%             diagrama con flechas revienta con «\language@active@arg> has an
%             extra }». Debe ir ANTES de que se dibuje nada.
%  · positioning        → sintaxis «below left=2pt and 3pt».
%  · decorations.pathreplacing → llaves ({brace}) para señalar rangos.
\usetikzlibrary{babel,positioning,decorations.pathreplacing}
\usepackage{graphicx}
% Circuitos: el trabajo de electrónica se hace hoy con micro:bit y sus sensores
% integrados (MakeCode, sin cableado), así que no se necesitan esquemáticos.
% Si algún día entran montajes con protoboard, basta descomentar esta línea:
% los diagramas .tex/.tikz declarados en `recursos:` ya se incrustan solos.
% \usepackage{circuitikz}
\usepackage[most]{tcolorbox}
\usepackage{tabularx}
\usepackage{array}
\usepackage{fancyhdr}
\usepackage{titlesec}
\usepackage{ragged2e}
\usepackage{enumitem}
\usepackage{microtype}

% Helvetica Neue no trae la flecha U+2192, asi que cada «→» escrito literalmente
% en el YAML se compilaba a NADA, en silencio (462 flechas en 61 guias: rutas del
% tipo «paleta → tipografia → lockup» salian rotas). Mapeamos el caracter a su
% version matematica, que si tiene glifo. Asi el autor puede seguir escribiendo
% «→» comodamente en el YAML.
\usepackage{newunicodechar}
\newunicodechar{→}{\ensuremath{\rightarrow}}

\setmainfont{Helvetica Neue}
\setsansfont{Helvetica Neue}
\setlength{\parindent}{0pt}
\setlength{\parskip}{6pt}
\hyphenpenalty=200
\exhyphenpenalty=200
\pretolerance=400
\tolerance=2000
\setlength{\emergencystretch}{1em}
\renewcommand{\arraystretch}{1.25}
\setlist{leftmargin=5mm,itemsep=1mm,topsep=1mm}
\newcolumntype{Y}{>{\RaggedRight\arraybackslash}X}

\definecolor{milcMagenta}{HTML}{E6007E}
\definecolor{milcVino}{HTML}{5A0038}
\definecolor{milcTurquesa}{HTML}{0093A5}
\definecolor{milcVerde}{HTML}{58A923}
\definecolor{milcMostaza}{HTML}{E5B400}
\definecolor{milcOcre}{HTML}{B86600}
\definecolor{milcNegro}{HTML}{241816}
\definecolor{milcGris}{HTML}{F7F7F4}
\definecolor{milcLinea}{HTML}{E8E2DD}
\definecolor{milcGradePrimary}{HTML}{84CC16}
\definecolor{milcGradeDark}{HTML}{4D7C0F}
\definecolor{milcGradeSoft}{HTML}{ECFCCB}
\definecolor{milcGradeAccent}{HTML}{CFE7FF}
\definecolor{milcGradeText}{HTML}{FFFFFF}
\color{milcNegro}

\pagestyle{fancy}
\fancyhf{}
\lhead{\small Tecnología e Informática · Bebras}
\rhead{\small Momento 9 · MILC v3}
\cfoot{\small\thepage}
\renewcommand{\headrulewidth}{0pt}

\newtcolorbox{titlebox}[1]{
  enhanced,colback=#1,colframe=#1,arc=7mm,boxrule=0pt,
  left=7mm,right=7mm,top=3mm,bottom=3mm,
  before skip=8pt,after skip=8pt,
  fontupper=\sffamily\bfseries\Large\color{white}
}

\newtcolorbox{softbox}[3]{
  enhanced,colback=#2,colframe=#1,borderline west={4pt}{0pt}{#1},
  arc=2mm,boxrule=.4pt,left=4mm,right=4mm,top=3mm,bottom=3mm,
  before skip=6pt,after skip=8pt,title={#3},coltitle=white,
  colbacktitle=#1,fonttitle=\sffamily\bfseries,fontupper=\RaggedRight,
  attach boxed title to top left={xshift=3mm,yshift=-2mm},
  boxed title style={arc=2mm, boxrule=0pt}
}

\newtcolorbox{drawbox}[2]{
  enhanced,colback=white,colframe=#1,arc=4mm,boxrule=1pt,
  left=5mm,right=5mm,top=4mm,bottom=4mm,before skip=8pt,after skip=8pt,
  title={#2},coltitle=white,colbacktitle=#1,fonttitle=\sffamily\bfseries,
  fontupper=\RaggedRight,
  attach boxed title to top left={xshift=4mm,yshift=-2mm},
  boxed title style={arc=3mm, boxrule=0pt}
}

\newtcolorbox{infoband}[1]{
  enhanced,colback=#1!8,colframe=#1!50,arc=2mm,boxrule=.5pt,
  left=4mm,right=4mm,top=2mm,bottom=2mm,before skip=4pt,after skip=4pt,
  fontupper=\small\RaggedRight
}

% Puente narrativo entre secciones (contrato editorial Conéctate).
% Da continuidad: "Acabas de X. Ahora vas a Y porque Z."
% Visualmente: banda izquierda en color del grado, texto en itálica pequeña.
\newtcolorbox{puentebox}{
  enhanced,colback=milcGradeSoft,colframe=milcGradeDark,
  borderline west={3pt}{0pt}{milcGradeDark},
  arc=0mm,boxrule=0pt,left=5mm,right=4mm,top=2.5mm,bottom=2.5mm,
  before skip=4pt,after skip=8pt,
  fontupper=\itshape\small\color{milcNegro}\RaggedRight
}
% La flecha va en modo matematico a proposito: Helvetica Neue NO tiene el glifo
% U+2192, asi que \textrightarrow se compilaba a NADA («Missing character: There
% is no → in font Helvetica Neue») y el puente salia sin flecha. En modo
% matematico la flecha viene de la fuente matematica y si se dibuja.
\newcommand{\puente}[1]{%
  \begin{puentebox}%
    \textcolor{milcGradeDark}{$\rightarrow$}\hspace{2mm}#1%
  \end{puentebox}%
}

\newcommand{\linead}[1]{\noindent\color{milcLinea}\rule{#1}{0.5pt}\color{milcNegro}}
\newcommand{\respuesta}[1]{\par\vspace{2mm}\foreach \n in {1,...,#1}{\noindent\color{milcLinea}\rule{\linewidth}{0.5pt}\par\vspace{5mm}}\color{milcNegro}}
\newcommand{\checkbox}{\textcolor{milcGradeDark}{\fbox{\phantom{x}}}\hspace{5pt}}

\newcommand{\phasecard}[4]{
\begin{tcolorbox}[enhanced,colback=#3,colframe=#2,arc=3mm,boxrule=.5pt,height=3.05cm,valign=center,halign=center,left=1.5mm,right=1.5mm,top=2mm,bottom=2mm]
{\sffamily\bfseries\fontsize{22}{24}\selectfont\textcolor{#2}{#1}}\par
{\sffamily\bfseries\small #4}
\end{tcolorbox}
}

% ── Recursos visuales de la guía ──────────────────────────────────────
% Declarados en el bloque `recursos:` del YAML y emitidos por el builder.
% \guiaFigura   → imagen rasterizada o PDF (foto, carta celeste, montaje).
% \guiaDiagrama → fuente TikZ/circuitikz (.tex) incrustada como vector:
%                 es la vía para circuitos y esquemáticos editables.
% #1 = ruta relativa al .tex · #2 = pie (puede ir vacío)
\newcommand{\guiaFigura}[2]{%
\begin{center}
\includegraphics[width=0.88\linewidth,height=0.38\textheight,keepaspectratio]{#1}\\[1.5mm]
{\footnotesize\itshape\textcolor{milcNegro!75}{#2}}
\end{center}
}

\newcommand{\guiaDiagrama}[2]{%
\begin{center}
\input{#1}\\[1.5mm]
{\footnotesize\itshape\textcolor{milcNegro!75}{#2}}
\end{center}
}

\begin{document}
\thispagestyle{empty}
\begin{tikzpicture}[remember picture,overlay]
  \fill[white] (current page.south west) rectangle (current page.north east);
  \path[fill=milcGradePrimary,rounded corners=16mm]
    ([xshift=.55cm,yshift=-.55cm]current page.north west) rectangle
    ([xshift=-.55cm,yshift=3.65cm]current page.south east);

  \node[anchor=north west,text=milcGradeText] at ([xshift=2.0cm,yshift=-1.85cm]current page.north west)
    {\sffamily\bfseries\fontsize{14}{16}\selectfont MOMENTO};
  \node[anchor=north west,text=milcGradeText,opacity=.82] at ([xshift=2.0cm,yshift=-2.75cm]current page.north west)
    {\sffamily\bfseries\fontsize{9}{11}\selectfont BEBRAS · PENSAR COMO UNA MÁQUINA};

  \begin{scope}[shift={([xshift=-3.05cm,yshift=-2.55cm]current page.north east)}]
    \fill[white,opacity=.80] (-.62,-.52) rectangle (.62,.52);
    \draw[milcGradeText,line width=1pt] (-.62,-.52) rectangle (.62,.52);
    \fill[milcGradeDark] (-.42,-.38) rectangle (-.22,.06);
    \fill[milcGradeAccent] (-.10,-.38) rectangle (.10,.24);
    \fill[milcGradePrimary!60!black] (.22,-.38) rectangle (.42,.38);
  \end{scope}

  \node[anchor=west,text=milcGradeText] at ([xshift=1.9cm,yshift=-12.85cm]current page.north west)
    {\sffamily\bfseries\fontsize{124}{124}\selectfont 9};
  % El circulito de grado (°) solo aplica a las guias de grado («11°»); en
  % Bebras y Semillero el numero grande es un momento/modulo, no un grado.
  

  \node[anchor=west,text=milcGradeText,text width=8.7cm,align=left] at ([xshift=10.35cm,yshift=-11.6cm]current page.north west)
    {
     {\sffamily\bfseries\fontsize{19}{22}\selectfont Cierre en minga --- simulacro integrador\\y diseño de tu propia tarea Bebras}\\[4mm]
     {\sffamily\bfseries\fontsize{23}{26}\selectfont Bebras · Momento 9}\\[2mm]
     {\sffamily\fontsize{13}{16}\selectfont Curso «Pensar como una máquina» · Pensamiento computacional}\\[5mm]
     {\sffamily\bfseries\fontsize{11}{13}\selectfont Institución Educativa Sor María Juliana}\\[1mm]
     {\sffamily\fontsize{10}{12}\selectfont Autor: Dr. Álvaro Cárdenas Orozco}
    };

  \path[fill=milcGradeSoft,rounded corners=7mm]
    ([xshift=1.35cm,yshift=.85cm]current page.south west) rectangle
    ([xshift=-1.35cm,yshift=4.25cm]current page.south east);
  \draw[milcGradeDark,line width=.65pt,rounded corners=7mm]
    ([xshift=1.35cm,yshift=.85cm]current page.south west) rectangle
    ([xshift=-1.35cm,yshift=4.25cm]current page.south east);
  \node[anchor=center,text=milcNegro,text width=17.0cm,align=center] at ([yshift=2.55cm]current page.south)
    {\sffamily\fontsize{10}{12}\selectfont Metodología MILC · Apertura ancestral · Pensamiento computacional · Triángulo Dussel-Estoicismo-Floridi\\
    \textcolor{milcGradeDark}{\bfseries Producto:} Tu tarea Bebras original (historia + 4 opciones + explicación del dominio) y la evaluación de la tarea de un compañero; más el simulacro integrador de 6 retos resuelto y autoevaluado.};
\end{tikzpicture}
\mbox{}
\newpage

\section*{Apertura: Cierre en minga: simulacro integrador y diseño de tu propia tarea Bebras}

\begin{softbox}{milcGradeDark}{milcGradeSoft}{Saber ancestral}
En las veredas del Valle del Cauca, la \textbf{minga} no termina cuando cae la tarde: termina cuando el trabajo está hecho \emph{y} cuando el grupo cierra en círculo, come el sancocho y cuenta lo que aprendió. Ese momento de cierre no es protocolo --- es parte del método. La comunidad no se dispersa sin antes preguntarse: ¿qué salió bien? ¿qué cambiaríamos la próxima vez? ¿quién necesita que alguien le explique lo que aquí aprendimos? Empezamos este curso con la minga como apertura --- la comunidad que parte una tarea enorme en partes pequeñas, que repite lo que funcionó, que se queda con lo esencial. Hoy la minga vuelve, pero esta vez como \textbf{cierre}: lo que aprendiste en nueve momentos --- descomponer, ver patrones, abstraer, seguir algoritmos, codificar datos, deducir con lógica, recorrer redes y competir con estrategia --- no es un tesoro para guardarlo solo. El pensamiento computacional es una herramienta de la comunidad.
\end{softbox}

\begin{softbox}{milcTurquesa}{milcGris}{Saber contemporáneo}
El \textbf{pensamiento computacional} tiene ocho dominios, y este curso los recorrió todos: descomposición, reconocimiento de patrones, abstracción, algoritmos, representación de datos, lógica y deducción, grafos y optimización, y estrategia de concurso. Cada dominio fue entrenado con tareas estilo \textbf{Bebras} --- retos de 45 minutos que se resuelven con la cabeza, no con código. El capstone de cualquier aprendizaje real no es resolver el examen final: es demostrar que \emph{puedes enseñar} lo que aprendiste. Por eso, la actividad cumbre de este momento es \textbf{crear tu propia tarea Bebras}: inventar una historia contextualizada en Cartago o en el Valle, formular una pregunta con cuatro opciones (una sola correcta), y explicar qué dominio del pensamiento computacional usa. Quien puede diseñar un buen reto, entendió de verdad.
\end{softbox}

\begin{softbox}{milcMostaza}{milcGris}{Pregunta-puente}
\textit{La minga siempre cierra con una pregunta: ¿para qué y para quién vas a usar lo que aprendiste? Nueve momentos de pensamiento computacional te dan una herramienta poderosa. ¿Qué harás con ella?}
\end{softbox}

\begin{softbox}{milcVerde}{milcGris}{Saber y saber hacer}
Al terminar podrás: (1) \textbf{integrar} los 8 dominios del pensamiento computacional en un simulacro cronometrado que replica las condiciones del concurso real; (2) \textbf{evaluar} el trabajo de un compañero con criterios explícitos y retroalimentación constructiva; (3) \textbf{crear} una tarea Bebras original con historia contextualizada, cuatro opciones plausibles, respuesta correcta justificada y dominio identificado; (4) \textbf{cerrar} el curso con un horizonte filosófico sobre para qué y para quién sirve aprender a pensar computacionalmente.
\end{softbox}



\small
\begin{tabularx}{\textwidth}{>{\bfseries}p{3.0cm}Y}
\rowcolor{milcGradePrimary!12}Autor & Dr. Álvaro Cárdenas Orozco\\
DBA & Desarrolla el pensamiento computacional --- descomposición, patrones, abstracción y algoritmos --- para resolver problemas (transversal 6°--11°, alineado a los lineamientos MEN de Tecnología e Informática)\\
Referentes & Valentina Dagienė (Bebras) · Pensamiento computacional (Wing) · Dussel · Estoicismo · Floridi · Saberes del Valle y el Pacífico\\
Producto final & Tu tarea Bebras original (historia + 4 opciones + explicación del dominio) y la evaluación de la tarea de un compañero; más el simulacro integrador de 6 retos resuelto y autoevaluado.\\
\end{tabularx}
\normalsize

\puente{La minga cierra en círculo. Mira el mapa de esta sesión: vas a repasar el camino completo, luego a crear --- no solo a resolver --- y finalmente a reflexionar sobre lo que cambió en ti.}

\begin{titlebox}{milcGradeDark}
Ruta MILC: así vamos a aprender
\end{titlebox}
\begin{tabularx}{\textwidth}{>{\centering\arraybackslash}X>{\centering\arraybackslash}X>{\centering\arraybackslash}X>{\centering\arraybackslash}X}
\phasecard{1}{milcVerde}{milcGris}{Escucha\\[-1mm]\small Observo, escucho y pregunto.} &
\phasecard{2}{milcTurquesa}{milcGris}{Sistematización\\[-1mm]\small Pensamiento computacional.} &
\phasecard{3}{milcMagenta}{milcGris}{Praxis\\[-1mm]\small Construyo y aplico.} &
\phasecard{4}{milcVino}{milcGris}{Evaluación\\[-1mm]\small Reflexión liberadora.}
\end{tabularx}


\puente{Antes de crear algo nuevo, vale la pena mirar atrás. Recuerda el primer momento del curso y pregúntate qué habría cambiado si ya supieras lo que sabes ahora.}

\begin{titlebox}{milcVerde}
Fase 1 · Escucha
\end{titlebox}

\begin{softbox}{milcVerde}{milcGris}{Escena para escuchar}
\textbf{Vuelve al momento 0}. Recuerda cuando llegaste a la primera sesión de este curso sin saber qué era Bebras ni qué era el pensamiento computacional. Piensa en un problema --- de tu casa, tu barrio o tu colegio --- que en ese entonces te parecía demasiado enredado. Ahora, con los ocho dominios en la cabeza, \emph{¿cómo lo partirías?} ¿Qué patrones verías? ¿Qué información descartarías? ¿Qué pasos seguirías? No tienes que resolver el problema entero: solo observa cómo cambió tu manera de mirarlo.
\end{softbox}

\checkbox Nombraste el problema concreto que recordaste --- no uno abstracto.\\
\checkbox Señalaste al menos dos dominios que usarías para abordarlo (por nombre, no solo por intuición).\\
\checkbox Escribiste algo que \emph{no harías} ahora que antes habrías intentado --- señal de que tu pensamiento cambió.

\begin{infoband}{milcVerde}


\textbf{En tu cuaderno (Actividad 1 \textbullet{} EXPLICA).} Título: «Actividad 1 --- El problema que ahora sí puedo partir». \textbf{Formato:} párrafo de 4 a 6 renglones. \textbf{Extensión:} describe el problema, di qué dominios usarías y qué cambió en tu forma de mirarlo. \textbf{Sabes que terminaste cuando} mencionas al menos dos dominios por nombre y explicas en qué te ayudan.
\end{infoband}

\puente{Acabas de mirar atrás. Ahora pon nombre a cada pieza del camino recorrido: los ocho dominios, uno a uno, para que el mapa completo quede claro antes del simulacro.}

\begin{titlebox}{milcTurquesa}
Fase 2 · Sistematización con pensamiento computacional
\end{titlebox}

\begin{softbox}{milcTurquesa}{milcGris}{Cuatro pilares aplicados al tema}
Nueve momentos. Ocho dominios. Cada uno fue una pieza del pensamiento computacional --- la misma forma de pensar que usa un computador y que tú ya usabas sin saberlo. Repasa el mapa completo antes del simulacro. No hay información nueva: todo lo que ves aquí ya lo trabajaste. Lo que cambia es que ahora lo ves junto, como un sistema.
\end{softbox}

\small
\begin{tabularx}{\textwidth}{>{\bfseries\RaggedRight\arraybackslash}p{3.6cm}Y}
\rowcolor{milcTurquesa!90}\color{white}Pilar & \color{white}\bfseries Aplicación al tema\\
1. Descomposición & \textbf{Las cuatro piezas fundacionales:} descomposición (partir lo grande en partes), reconocimiento de patrones (ver lo que se repite), abstracción (quedarse con lo esencial) y algoritmos (pasos exactos en el orden correcto). Son la base de todo lo demás. Sin ellas, los otros cuatro dominios no tienen dónde apoyarse.\\
2. Reconocimiento de patrones & \textbf{Representación de datos y lógica:} codificar información para que una máquina la procese (binario, tablas, grafos) y deducir la verdad descartando lo imposible (contrapositiva, tablas de verdad). Estos dos dominios convierten el pensamiento computacional en un lenguaje preciso.\\
3. Abstracción & \textbf{Grafos y optimización:} modelar redes y encontrar el camino de menor costo. Los ríos del Pacífico, las rutas del bus de Cartago y el internet son grafos: hay caminos mejores que otros, y saberlos encontrar es una ventaja real.\\
4. Algoritmo conceptual & \textbf{Estrategia de concurso:} saber cuándo arriesgar y cuándo guardar. En Bebras, un reto nivel C con penalización de $-4$ puntos no se ataca a ciegas. La estrategia es parte del pensamiento computacional: no toda la información disponible merece ser procesada.\\
\end{tabularx}
\normalsize

\begin{drawbox}{milcTurquesa}{Mapa de los 8 dominios: cómo se encadenan}
\textbf{1. Descomposición} --- partir el problema. \\[1mm] \textbf{2. Patrones} --- ver lo que se repite. \\[1mm] \textbf{3. Abstracción} --- quedarse con lo esencial. \\[1mm] \textbf{4. Algoritmos} --- pasos exactos en orden. \\[1mm] \textbf{5. Datos} --- codificar para que la máquina entienda. \\[1mm] \textbf{6. Lógica} --- deducir descartando lo imposible. \\[1mm] \textbf{7. Grafos} --- recorrer redes, optimizar rutas. \\[1mm] \textbf{8. Estrategia} --- decidir cuándo actuar y cuándo esperar. \\[2mm] Los cuatro primeros son las \emph{piezas fundacionales}. Los cuatro siguientes los \emph{aplican a dominios concretos}. Un buen pensador computacional usa los ocho, no solo los que le gustan.
\end{drawbox}

\begin{softbox}{milcMostaza}{milcGris}{Errores comunes y su corrección}
\textbf{Creer que los ocho dominios son independientes:} descomponer sin abstraer lleva a partes innecesarias. Abstraer sin algoritmo no llega a ningún lado. Los dominios se usan juntos. \textbf{Confundir lógica con sentido común:} la lógica formal es más estricta --- «me parece que» no es una deducción válida. \textbf{Ignorar la estrategia:} en el concurso real, un error en nivel C cuesta más que no responder. Calcular el riesgo \emph{es} pensar computacionalmente.
\end{softbox}

\begin{infoband}{milcTurquesa}


\textbf{En tu cuaderno (Actividad 2 \textbullet{} ANALIZA).} Título: «Actividad 2 --- Los 8 dominios en un solo mapa». \textbf{Formato:} tabla de 2 columnas (dominio · ejemplo de la vida cotidiana en Cartago o el Valle). \textbf{Extensión:} los ocho dominios, cada uno con un ejemplo distinto a los del texto. \textbf{Sabes que terminaste cuando} tus ejemplos no repiten los del curso y un compañero los reconocería sin que se los expliques.
\end{infoband}

\puente{Ya tienes el mapa completo. Es hora de demostrar lo que sabes --- no resolviendo un reto más, sino \textbf{creando} el tuyo: una tarea Bebras original, contextualizada en Cartago, que otro pueda resolver.}

\begin{titlebox}{milcMagenta}
Fase 3 · Praxis con pensamiento computacional
\end{titlebox}

\begin{softbox}{milcMagenta}{milcGris}{Construyendo el producto con los cuatro pilares}
El dominio más alto del aprendizaje no es resolver: es \textbf{crear}. Vas a diseñar tu propia tarea estilo Bebras --- una historia corta ambientada en Cartago, el Valle del Cauca o en tu barrio, con una pregunta de cuatro opciones (una sola correcta), y la explicación de qué dominio del pensamiento computacional usa. Quien puede diseñar un buen reto, entendió de verdad. Después, vas a evaluar la tarea de un compañero con criterios claros.
\end{softbox}

\small
\begin{tabularx}{\textwidth}{>{\bfseries\RaggedRight\arraybackslash}p{3.6cm}Y}
\rowcolor{milcMagenta!90}\color{white}Pilar & \color{white}\bfseries Aplicación al producto\\
Descomposición & \textbf{Elige el dominio:} uno de los ocho del curso. No elijas el más fácil para ti: elige el que más te interesa enseñar. El dominio guía todo lo demás --- la historia, la pregunta y la explicación.\\
Patrones & \textbf{Construye la historia:} ambientada en Cartago, el Valle o tu barrio. Debe ser corta (3 a 5 renglones), concreta y reconocible para tus compañeros. Evita nombres de personajes genéricos como «Juan» --- usa lugares reales: la galería de Cartago, el río La Vieja, el cerro El Bolo.\\
Abstracción & \textbf{Formula la pregunta con 4 opciones:} una sola respuesta correcta, las otras tres \emph{plausibles pero descartables}. Las opciones incorrectas son el corazón de un buen reto Bebras: deben parecer posibles para quien no pensó con cuidado, pero imposibles para quien sí usó el dominio correcto.\\
Algoritmo de armado & \textbf{Escribe la explicación:} di por qué la respuesta correcta es correcta, por qué las otras tres no sirven, y nombra el dominio del pensamiento computacional que usaste. Esta explicación demuestra que tú entendiste --- no que copiaste.\\
\end{tabularx}
\normalsize

\begin{drawbox}{milcMagenta}{Plantilla para diseñar tu tarea Bebras}
\checkbox \textbf{Dominio elegido:} \underline{\hspace{6cm}} \\[3mm]
\checkbox \textbf{Historia} (3--5 renglones, ambientada en Cartago, el Valle o tu barrio): \\[6mm] \underline{\hspace{\linewidth}} \\[2mm] \underline{\hspace{\linewidth}} \\[2mm] \underline{\hspace{\linewidth}} \\[3mm]
\checkbox \textbf{Pregunta:} \\[5mm] \underline{\hspace{\linewidth}} \\[3mm]
\checkbox \textbf{Opciones:} \quad A. \underline{\hspace{3.5cm}} \quad B. \underline{\hspace{3.5cm}} \quad C. \underline{\hspace{3.5cm}} \quad D. \underline{\hspace{3.5cm}} \\[3mm]
\checkbox \textbf{Respuesta correcta:} \quad \fbox{\phantom{x}}~A \quad \fbox{\phantom{x}}~B \quad \fbox{\phantom{x}}~C \quad \fbox{\phantom{x}}~D \\[3mm]
\checkbox \textbf{Explicación} (por qué las otras tres no sirven + dominio usado): \\[6mm] \underline{\hspace{\linewidth}} \\[2mm] \underline{\hspace{\linewidth}}
\end{drawbox}

\begin{softbox}{milcMagenta}{milcGris}{Plantilla de trabajo}
\textbf{Evaluación entre pares (Actividad 4 \textbullet{} EVALÚA).} Intercambia tu cuaderno con el de un compañero. Lee su tarea sin mirar la explicación. Intenta resolver la pregunta solo con las cuatro opciones. Luego evalúa con tres criterios: (1) ¿la historia es clara y está bien contextualizada? (2) ¿hay una única respuesta posible, o más de una opción podría ser correcta? (3) ¿el dominio que señala es el correcto? Escribe tu veredicto --- no «está bien» sino razones concretas --- y una sugerencia de mejora específica y respetuosa. \textbf{En tu cuaderno (Actividad 4):} Título: «Actividad 4 --- Evaluación de la tarea de mi compañero». \textbf{Formato:} tres viñetas (una por criterio) + 2 a 3 renglones de sugerencia. \textbf{Sabes que terminaste cuando} tu veredicto en cada criterio es concreto y tu sugerencia señala algo específico que el compañero puede mejorar.
\end{softbox}

\begin{infoband}{milcMagenta}


\textbf{En tu cuaderno (Actividad 3 \textbullet{} CREA).} Título: «Actividad 3 --- Mi tarea Bebras». \textbf{Formato:} historia de 3 a 5 renglones + pregunta con 4 opciones claramente marcadas (A, B, C, D) + explicación de 3 a 5 renglones. \textbf{Extensión:} mínimo una página de cuaderno. \textbf{Sabes que terminaste cuando} tu historia está ambientada en un contexto reconocible, hay una sola respuesta posible, y tu explicación nombra el dominio del pensamiento computacional y justifica por qué las otras tres opciones no sirven.
\end{infoband}

\puente{Diseñaste tu tarea y evaluaste la de un compañero. Ahora pon a prueba los ocho dominios en el simulacro: seis retos, uno por dominio, en silencio, como en el concurso real.}

\begin{titlebox}{milcMostaza}
Producto de la sesión
\end{titlebox}
\textbf{Tu tarea Bebras original y la evaluación de la tarea de tu compañero}.
\begin{softbox}{milcMostaza}{milcGris}{Criterios de calidad}
(1) La historia de tu tarea está ambientada en Cartago, el Valle del Cauca o en un contexto cotidiano que tus compañeros reconocen. (2) Hay una única respuesta correcta entre las cuatro opciones; las otras tres son plausibles pero descartables con el dominio correcto. (3) Tu explicación nombra el dominio del pensamiento computacional y justifica por qué las otras tres opciones son incorrectas. (4) Tu evaluación de la tarea del compañero tiene un veredicto concreto por cada uno de los tres criterios --- no «está bien» sino razones específicas. (5) En el simulacro integrador, identificaste en qué dominios te sientes sólido y en cuáles todavía hay trabajo, sin importar el puntaje.
\end{softbox}

\puente{Terminaste el simulacro. Detente un momento: ¿en qué dominios te sentiste sólido? ¿En cuáles todavía hay trabajo por hacer?}

\begin{titlebox}{milcVino}
Fase 4 · Evaluación liberadora — cinco dimensiones
\end{titlebox}

\begin{softbox}{milcVino}{milcGris}{Sentido de la evaluación}
Evaluar no es solo revisar una nota. Es reconocer qué aprendí, qué cambió en mí, qué emociones manejé, cómo me relaciono con la ciudad y la comunidad, y qué le dejaré a quien viene detrás.
\end{softbox}

\textbf{Mi reflexión liberadora en cinco dimensiones.} Completa cada frase iniciada:

\small
\begin{tabularx}{\textwidth}{>{\bfseries\RaggedRight\arraybackslash}p{3.4cm}Y>{\RaggedRight\arraybackslash}p{4.6cm}}
\rowcolor{milcVino}\color{white}Dimensión & \color{white}\bfseries Mi respuesta (frame starter) & \color{white}\bfseries Algo que descubrí\\
1. Personal & Lo que más cambió en mí fue \linead{6.5cm} & \linead{4.2cm}\\
2. Emocional & La emoción más fuerte fue \linead{2.5cm} y la regulé \linead{3cm} & \linead{4.2cm}\\
3. Ciudadana & En democracia esto importa porque \linead{6cm} & \linead{4.2cm}\\
4. Local (Cartago) & En mi barrio sería útil para \linead{6.5cm} & \linead{4.2cm}\\
5. Intergeneracional & A mi abuela/abuelo le diría \linead{6.5cm} & \linead{4.2cm}\\
\end{tabularx}
\normalsize

\begin{infoband}{milcVino}
\textbf{Para la próxima sustentación.} Vuelve a esta tabla en seis meses. Verás que las respuestas cambian — no porque lo aprendido cambie, sino porque tú cambias. La evaluación liberadora es una conversación contigo a través del tiempo.
\end{infoband}


\puente{Cerramos el curso pensando en grande. Tres voces filosóficas para entender qué significa --- para ti y para tu comunidad --- haber aprendido a pensar como una máquina.}

\begin{titlebox}{milcGradeDark}
Triángulo de pensamiento · cierre filosófico
\end{titlebox}

\begin{softbox}{milcVino}{milcGris}{Dussel · lente del nosotros}
\textit{``El otro ---el pobre, el excluido--- es el punto de partida de toda ética. Pensar bien no es un privilegio: es una herramienta de liberación para quienes más la necesitan.''}

La minga nunca es para uno solo. Aprender a descomponer, ver patrones y deducir con lógica es una ventaja --- y las ventajas, en una ética del nosotros, se comparten. Bebras no era el objetivo: era el entrenamiento para pensar con más claridad al servicio de tu comunidad.

\textbf{¿Quién en tu barrio, tu familia o tu colegio necesita que alguien con estas herramientas le ayude a resolver un problema? ¿Cómo podrías ser esa persona?}
\end{softbox}
\linead{\linewidth}\\[1mm]
\linead{\linewidth}

\begin{softbox}{milcMostaza}{milcGris}{Estoicismo · Marco Aurelio · cuidado interior}
\textit{``El conocimiento que no mejora al que lo tiene es inútil. Aprende no para acumular, sino para vivir con más claridad.''}

Nueve momentos, ocho dominios, decenas de retos. Todo ese recorrido solo tiene valor si cambió algo en la forma en que enfrentas un problema cuando nadie te está mirando. No la nota, no el concurso: la claridad con que piensas cuando estás solo frente a algo difícil.

\textbf{¿En quién te convierte haber aprendido a pensar así? ¿Hay algo en tu forma de enfrentar los problemas que ya cambió, aunque no te hayas dado cuenta?}
\end{softbox}
\linead{\linewidth}\\[1mm]
\linead{\linewidth}

\begin{softbox}{milcTurquesa}{milcGris}{Floridi · lente de la infoesfera}
\textit{``Somos inforgs: habitamos y somos constituidos por la infoesfera. Quien la entiende tiene una responsabilidad que quien no la entiende no puede tener.''}

Ahora sabes cómo piensa una máquina. Sabes cómo codifica datos, cómo recorre redes, cómo deduce con lógica. Eso no es solo conocimiento técnico: es poder. Y el poder, en la infoesfera que ya habitamos, viene con una responsabilidad que antes no tenías.

\textbf{Ahora que sabes cómo piensa una máquina, ¿qué harás diferente en la forma en que usas la tecnología, consumes información o resuelves problemas en tu comunidad?}
\end{softbox}
\linead{\linewidth}\\[1mm]
\linead{\linewidth}

\begin{titlebox}{milcVino}
Compromiso final
\end{titlebox}

\small
\begin{tabularx}{\textwidth}{>{\RaggedRight\arraybackslash}p{3.0cm}YYY}
\rowcolor{milcVino}\color{white}\bfseries Criterio & \color{white}\bfseries Lo logré & \color{white}\bfseries En proceso & \color{white}\bfseries Necesito apoyo\\
Comprensión del tema & Explico ideas centrales con mis palabras. & Reconozco ideas pero falta precisión. & Necesito repasar conceptos base.\\
Pensamiento computacional & Apliqué los 4 pilares al diseñar mi producto. & Apliqué algunos pilares. & Necesito releer la fase 2.\\
Aplicación y producto & Mi producto cumple los criterios y se puede revisar. & Producto funcional con ajustes pendientes. & Aún en construcción.\\
Reflexión 5 dimensiones & Respondí cada dimensión con honestidad. & Respondí algunas con detalle. & Mis respuestas son muy generales.\\
Triángulo de pensamiento & Conecté las 3 voces con mi trabajo real. & Conecté al menos una. & No relaciono las voces con mi proceso.\\
\end{tabularx}
\normalsize

\textbf{Hoy entendí que:}\\[1mm]
\linead{\linewidth}\\[1mm]
\linead{\linewidth}

\textbf{Mi compromiso para la próxima sesión es:}\\[1mm]
\linead{\linewidth}\\[1mm]
\linead{\linewidth}

\begin{softbox}{milcMostaza}{milcGris}{Cita de despedida · Marco Aurelio}
\textit{``El obstáculo en el camino se vuelve el camino. Lo que estorba al camino se vuelve camino.''}\\[1mm]
Cada error de hoy, bien observado, se convierte en pieza del camino que pisarás mañana.
\end{softbox}

\small
\begin{tabularx}{\textwidth}{>{\bfseries}p{3.6cm}Y}
\rowcolor{milcGradeDark!12}Nombre completo: & \linead{10cm}\\
Fecha: & \linead{4cm} \quad Curso/grupo: \linead{4cm}\\
Mi firma: & \linead{9cm}\\
\end{tabularx}
\normalsize

\begin{infoband}{milcGradeDark}
\textbf{Cierre del curso.} El pensamiento computacional no termina aquí: lo llevas contigo. Esta semana, cuando enfrentes un problema difícil --- en matemáticas, en casa, en el barrio ---, detente y pregúntate: ¿cómo lo partiría? ¿qué se repite? ¿qué puedo ignorar? ¿qué pasos debo seguir? Si quieres seguir practicando con retos reales, entra al archivo oficial en \texttt{bebras.org} y resuelve las tareas de años anteriores. Si quieres ir más lejos, habla con tu docente sobre participar en el concurso real. La minga terminó --- pero el pensamiento que cultivamos aquí, no.
\end{infoband}

\end{document}
